\chapter{Named Entity Recognition and
Summarizing}\label{sec-named-entity-recognition-and-summarization}

\textbf{Learning Objectives}

\emph{Overview starts}

\begin{itemize}
\tightlist
\item
  Understanding what named entity recognition (NER) is.
\item
  Understanding how NER works.
\item
  Understanding workflow in NER.
\item
  Learning how to apply NER in R.
\item
  Learning how to extract information using spacyr().
\item
  Understanding and selecting parts-of-speech (POS) tagging.
\item
  Analyzing POS over a timeline.
\item
  Measuring language complexity.
\item
  Tackling summarizing.
\end{itemize}

\emph{Overview ends}

Named entity recognition (NER) is used in NLP to find important
information in documents. Its goal is to identify and classify basic
text elements and their interactions. Named entities (NE) are important
for a specific purpose or topic, determining their type and boundary.
Although NEs do not have a direct and precise definition, maybe because
there are too many of them, they can be defined from a more practical point of
view. They are elements which describe categories, such as people,
organizations, locations, dates, times, percentages, and monetary
amounts, among others. For example, a person can be identified by a name
or an alias, and a location could be a city, village, or forest. These
two categories, person and location, might be related by the place of
birth of someone identified by its name. NER is a branch of NLP and helps us understand text questions such as ''who?'', ''what?'', ''where?'', and ''when?'',
through simple principles (Nouvel, Ehrmann, and Rosset 2016).

As we have already demonstrated in the previous chapter when tracking the
evolution of brand sentiment over time, in marketing this information
can be used to monitor several aspects of a product or brand. A company might want to know the
opinion of their customers (''who''), what they think about their products or
brands (''what''), in which region or country they are located (''where''), and at which point in
time they bought the product (''when''). They might also want to check on competitors for their development of technologies, national or international marketing strategies, or the pain
points of a brand's customer perceptions. Some reviewing platforms,
ranging in their focus from general topics (e.g., Yelp) to more
specialized areas, such as traveling (e.g., TripAdvisor) or movie review
platforms (e.g., IMDb), aim to collect information on products and link
them to consumer feedback (Popescu and Etzioni 2007). If well
extracted, NEs can be considered the building blocks of a summary,
giving insights into the information contained in documents (Hassel
2003).

An important application of NER is the possibility to anonymize the
personal data contained in a document; this is very important as data privacy is paramount. Personal information can be
hidden while the related data and the possibility of interpreting them can
be used for statistical purposes or the development of segment-specific
marketing campaigns.

\section{How does NER work?}\label{how-does-ner-work}

There are several ways of using NER, such as through rule-based systems,
statistical models, and neural-network models. Rule-based systems use
predefined lists of entities, for example, gazetteers,\footnote{A
  gazetteer is akin to a geographical dictionary or directory} and
regular expressions, like the ones we presented in
Chapter~\ref{sec-textual-data}. This is relatively simple to set up,
but it is not very scalable or flexible if you have large heterogeneous sets of 
data. Statistical methods and machine-learning algorithms use
methods, like those we discussed for classification in
Chapter~\ref{sec-classification}. Some other frequently used approaches
include hidden Markov models (HMM; Morwal, Jahan, and Chopra 2012),
conditional random fields (CRF; Sun et al. 2006), and maximum entropy
(ME) methods (Chieu and Ng 2003). These are more accurate than
the rule-based systems, but the data used for training must be of a
substantial size. The most accurate methods are the neural-network
models. They use techniques such as recurrent neural networks (RNN; Li
et al. 2015), convolutional neural networks (CNN; Zhu, Wang, and
Karlsson 2019), and transformers (Arkhipov et al. 2019). The computational power required by such models is high, and
they also need large annotated data sets. Even then, however, results are not
always perfect, and the process of identifying and classifying can be
faulty, especially with entities that are ambiguous or
context-dependent. Take, as an example, Amazon, which can be connected with
a forest (location) or a company (organization). Disambiguation is a
complex task, but most recent transformer technologies understand
context well and have made this more and more precise. It is necessary to
mention that variability introduces misclassification, for example, in
cases where an entity is referred to in different ways, like CR7 or
Cristiano Ronaldo: This task is called co-reference resolution, sometimes
called named entity resolution, and is extremely useful in text analytics.

\subsection{Workflow}\label{workflow}

In an NER task, raw text must undergo several transformations, from
pre-processing to the final NER output. Here, we will explain a generic
NER method. Pre-processing helps transform the text into a more
suitable format for the subsequent NER procedures. It starts by segmenting
the text into sentences, which are represented by lists of strings. This
is followed by tokenization, which helps to break down text to the desired level of detail and becomes a list of strings. Next is the
parts-of-speech (POS) tagging, which identifies which parts of speech
specific words belong to: Tokens can be nouns, verbs, adjectives, and so
on, and are very useful for entity recognition. In modern contexts,
tokens can be transformed into word embedding; we will discuss this in
Chapter~\ref{sec-extracting-information-from-embeddings}. Next comes the
entity detection, which helps classify the tokens into categories
using one of the models discussed above. During this step, the model identifies the boundaries of the
entity and its type. Then, relation detection comes into
action, with the final result being in the form of text tuples associated with
entities and relationships. See Figure~\ref{fig-ner}, inspired by Dias
et al. (2020), for a visual representation of the process.

\begin{figure}

\centering{

\includegraphics[width=0.7\textwidth,height=\textheight]{figure/chapter9/ner.png}

}

\caption{\label{fig-ner}Named Entity Recognition Task}

\end{figure}%

The landscape of linguistic annotations is made more complex by the use
of various formats, two of them being more common than others: conference of natural
language learning universal dependencies (CoNNL-U; cf. Project 2022; Schuster and Manning 2016) and text interchange format (TIF; Mullen
et al. 2018). The former represents a highly standardized and widespread
format for dependency structures across various languages. Its form is a
plain text file containing several fields for each token, and it becomes
useful when universal dependencies are needed. The latter,
used in R especially, offers a flexible structure for text package objects, such as objects from
corpora and tokens. All R packages which support this standard will
be able to interact between them nicely; two R packages are
\texttt{spacyr} and \texttt{quanteda}.

The choice of model depends on several factors, for example, the
dimension of the data set, the presence of labeled data, the computational
capacity available, or the degree of precision and accuracy needed from the results.

\section{NER using R}\label{ner-using-r}

There are several packages which can be used in R. Each package might use
a different model to annotate the text. One of them is
\texttt{nametagger}, another is \texttt{OpenNLP}, which is based on the
Apache OpenNLP library. Related to the latter is \texttt{entity},
which is a wrapper for \texttt{NLP} and \texttt{OpenNLP}. Another is
\texttt{crfsuite}, providing an interface with CRF. Yet another
is \texttt{udpipe}, a package developed by the Belgium Network of Open
Source Analytical Consultants (BNOSAC). Packages like
\texttt{nametagger}, \texttt{crfsuite}, or \texttt{udpipe} use CoNLL-U
models \footnote{CoNLL-U format:
  https://universaldependencies.org/format.html} while \texttt{OpenNLP}
has its own annotator. Last but not least, an important package,
\texttt{spacyr} (Benoit and Matsuo 2023), is used in R through an
interface with the Python \texttt{SpaCy} library (Vasiliev 2020).
\texttt{SpaCy} uses neural-network models and tries to predict entities
based on them. The models are trained using large data sets.
Depending on the need, they are optimized for speed, sm (small), for a
balance between speed and accuracy, md (medium), and for higher accuracy
but with more computational power needed, lg (large). These models exist
for several languages, such as English, German, French, and many others. A
pipeline for how \texttt{Spacy} carries out annotations is provided in
Figure~\ref{fig-pipeline}. The first component, ``tok2vec'', transforms
tokens into embedded vectors (there is more about this in
Chapter~\ref{sec-extracting-information-from-embeddings}) which represent
semantic meaning in a numerical space. These vectors are used by both
the ``tagger'', which assigns POS tags, and the ``parser'' which analyzes
the grammatical structure of the sentence. The ``attribute\_ruler''
applies rules to elements such as lemma, POS, etc., and helps the
``lemmatizer'' standardize words as their root form. The last
component, ``ner'', identifies named entities. It might loop back,
however, to ``tok2vec'' for better detection.

\begin{figure}

\centering{

\includegraphics[width=0.9\textwidth,height=\textheight]{figure/chapter9/pipeline-design.png}

}

\caption{\label{fig-pipeline}SpaCy pipeline design \newline (Source:
spaCy website)}

\end{figure}%

The entities (or categories) for NER that \texttt{SpaCy} extracts are basic; some names are intuitive, while others might need an
explanation. The examples below are not exhaustive.

\begin{itemize}
\tightlist
\item
  PERSON: People, including fictional, aliases, former names
\item
  NORP (Nationalities, or Religious or Political): Nationalities or
  religious or political groups
\item
  FAC (Facility): Buildings, airports, highways, bridges, hospitals
\item
  ORG: Companies, agencies, institutions, NGOs
\item
  GPE (Geopolitical Entity): Countries, states, cities, rural areas
\item
  LOC: Non-GPE locations, mountain ranges, bodies of water, deserts
\item
  PRODUCT: Objects, vehicles, food (not services.)
\item
  EVENT: Named hurricanes, battles, wars, sports events, political
  events, cultural events, scientific events, economic events
\item
  WORK\_OF\_ART: Titles of books, songs, paintings
\item
  LAW: Named documents made into laws
\item
  LANGUAGE: Any named language
\end{itemize}

SpaCy includes some extended categories, too:

\begin{itemize}
\tightlist
\item
  DATE: Absolute or relative dates or periods
\item
  TIME: Times smaller than a day
\item
  PERCENT: Percentage, including ``\%''
\item
  MONEY: Monetary values, including units
\item
  QUANTITY: Measurements, e.g., of weight or distance
\item
  ORDINAL: ``first'', ``second'', etc.
\item
  CARDINAL: Numerals, such as 1, 10, 100, three, seventeen
\end{itemize}

The previous categories might have a letter assigned after them, for
example, DATE\_B. This is because it follows the BILUO tagging scheme,
which stands for:

\begin{itemize}
\tightlist
\item
  B: Beginning of an entity
\item
  I: Inside of an entity (continuation of an entity started with B)
\item
  L: Last token of an entity
\item
  U: Unit entity (a single-token entity)
\item
  O: Outside of any entity
\end{itemize}

This is used when the entity spans multiple tokens, and the current
token is the first, inside, the last, etc., token in that sequence.

\section{NER Information extraction using
spacyr}\label{ner-information-extraction-with-spacyr}

In this subsection we illustrate the functionalities of the NER process by using the R package \texttt{spacyr}, along with
the sentiment dataset we already familiar with from
Chapter~\ref{sec-sentiment-analysis}. We will extract the most frequent
entities for two brands and compare their counts. Then we will
focus on the negative reviews and check how often specific entities are
mentioned. Adjectives are often related to emotions, and checking them gives us
more insights into how brands are perceived. However, adjectives can be related to
other words, so we will also explore the part-of-speech (POS) dependencies.
Tracking adjectives over time is a very useful way to measure the levels of
satisfaction or dissatisfaction with products. Language complexity is
another metric which can be calculated, and is potentially important for
understanding and comparing how people express their opinions about
specific brands, and also gives a hint to managers for how to respond and match
the users' styles.

\subsection{Identifying and describing
entities}\label{identifying-and-describing-entities}

After loading the \texttt{spacyr} library, the package
introduces a function called \texttt{spacy\_initialize()} to open the
connection between R and Python, and which detects if the library exists and, if so,
initializes it. Frequently, the cohabitation of R and Python can be
problematic due to the environments Python might exist in and the versions or
way it was installed on your operating system. For this
data set, we suggest installing the model \texttt{en\_core\_web\_lg.}

We will use the \texttt{comparison\ dataset} from
Chapter~\ref{sec-sentiment-analysis}, which is a reduced version of the
initial data set of product reviews. It contains the sentiment scores
already assigned to each document, and we will select those which have a
sentiment.ai score \(<-0,7\), an empirical threshold, so we can focus on the
most negative reviews to identify the problems with
the brands. We have also dropped some fields which we will not need, such as
sentimentAnalysis, vader, sentimentr, or syuzhet, and selected two brands.
In the example below we chose Samsung and Apple.

\begin{Shaded}
\begin{Highlighting}[]
\FunctionTok{library}\NormalTok{(spacyr)}
\FunctionTok{library}\NormalTok{(dplyr)}
\FunctionTok{library}\NormalTok{(lubridate)}

\CommentTok{\# Initialize spaCy with a large English model}
\FunctionTok{spacy\_initialize}\NormalTok{(}\AttributeTok{model =} \StringTok{"en\_core\_web\_lg"}\NormalTok{)}


\CommentTok{\# Clean and filter data, then add row numbers as Review\_id}
\NormalTok{data }\OtherTok{\textless{}{-}}\NormalTok{ comparison }\SpecialCharTok{\%\textgreater{}\%}
  \FunctionTok{select}\NormalTok{(}\SpecialCharTok{{-}}\NormalTok{sentimentAnalysis, }\SpecialCharTok{{-}}\NormalTok{sentimentr, }\SpecialCharTok{{-}}\NormalTok{vader, }\SpecialCharTok{{-}}\NormalTok{syuzhet) }\SpecialCharTok{\%\textgreater{}\%}
  \FunctionTok{filter}\NormalTok{(sentiment.ai }\SpecialCharTok{\textless{}} \SpecialCharTok{{-}}\FloatTok{0.7}\NormalTok{, Brand }\SpecialCharTok{\%in\%} \FunctionTok{c}\NormalTok{(}\StringTok{"Samsung"}\NormalTok{, }\StringTok{"Apple"}\NormalTok{)) }\SpecialCharTok{\%\textgreater{}\%}
  \FunctionTok{mutate}\NormalTok{(}\AttributeTok{Review\_id =} \FunctionTok{row\_number}\NormalTok{())}
\end{Highlighting}
\end{Shaded}

Next we use the \texttt{spacy\_parse()} function in the text column to
perform some initial linguistic analysis. The parsed text is summarized
as a table using the \texttt{parsedtxt} variable which contains the text
identifier (the number of the review), the sentence's position in the
text, the token's position in the sentence, the actual word and its base
(root) form, the grammatical category (e.g., noun, verb, adjective), and
the NER labels as classifications for organizations, products,
and so on.

\begin{Shaded}
\begin{Highlighting}[]
\CommentTok{\# Parse text to extract syntactic dependencies}
\NormalTok{parsedtxt }\OtherTok{\textless{}{-}} \FunctionTok{spacy\_parse}\NormalTok{(data}\SpecialCharTok{$}\NormalTok{target)}
\FunctionTok{head}\NormalTok{(parsedtxt)}
\end{Highlighting}
\end{Shaded}

By carrying out this operation, we have lost some information from the original
data set, but if we want to have a correspondence between each document
we still can merge this new data frame with an enhanced
version of the original data set. Because it is helpful for
the subsequent processing to separate the text into sentences, we transform
each review's identification in the \texttt{data} frame by adding
``text'' to the numerical value of the ``doc\_id'' field. After
performing this operation, we save the data set to a new variable called
\texttt{sent\_df}. We then merge this with the original data set to
include the Brand column. Finally, we combine this data set with the
\texttt{parsedtxt} data frame to create a new structure ready for
analysis.

\begin{Shaded}
\begin{Highlighting}[]
\FunctionTok{library}\NormalTok{(dplyr)}

\CommentTok{\# Create doc\_id, remove Review\_id}
\NormalTok{data }\OtherTok{\textless{}{-}}\NormalTok{ data }\SpecialCharTok{\%\textgreater{}\%}
  \FunctionTok{mutate}\NormalTok{(}\AttributeTok{doc\_id =} \FunctionTok{paste0}\NormalTok{(}\StringTok{"text"}\NormalTok{, }\FunctionTok{as.character}\NormalTok{(Review\_id))) }\SpecialCharTok{\%\textgreater{}\%}
  \FunctionTok{select}\NormalTok{(}\SpecialCharTok{{-}}\NormalTok{Review\_id)}

\CommentTok{\# Tokenize text into sentences and assign sentence IDs}
\NormalTok{sent\_df}\OtherTok{\textless{}{-}} \FunctionTok{spacy\_tokenize}\NormalTok{(data}\SpecialCharTok{$}\NormalTok{target, }\AttributeTok{what =} \StringTok{"sentence"}\NormalTok{, }
                         \AttributeTok{output =} \StringTok{"data.frame"}\NormalTok{, }
                         \AttributeTok{doc\_id =}\NormalTok{ data}\SpecialCharTok{$}\NormalTok{doc\_id) }\SpecialCharTok{\%\textgreater{}\%}
  \FunctionTok{rename}\NormalTok{(}\AttributeTok{sentence =}\NormalTok{ token) }\SpecialCharTok{\%\textgreater{}\%}
  \FunctionTok{group\_by}\NormalTok{(doc\_id) }\SpecialCharTok{\%\textgreater{}\%}
  \FunctionTok{mutate}\NormalTok{(}\AttributeTok{sentence\_id =} \FunctionTok{row\_number}\NormalTok{()) }\SpecialCharTok{\%\textgreater{}\%}
  \FunctionTok{ungroup}\NormalTok{()}


\CommentTok{\# Merge back with original data to include the Brand and R\_date}
\NormalTok{sent\_df }\OtherTok{\textless{}{-}}\NormalTok{ sent\_df }\SpecialCharTok{\%\textgreater{}\%}
  \FunctionTok{left\_join}\NormalTok{(data }\SpecialCharTok{\%\textgreater{}\%} \FunctionTok{select}\NormalTok{(doc\_id, Brand, R\_date), }\AttributeTok{by =} \StringTok{"doc\_id"}\NormalTok{ )}

\NormalTok{sentence\_data }\OtherTok{\textless{}{-}}\NormalTok{ parsedtxt }\SpecialCharTok{\%\textgreater{}\%}
  \FunctionTok{left\_join}\NormalTok{(sent\_df , }\AttributeTok{by =} \FunctionTok{c}\NormalTok{(}\StringTok{"doc\_id"}\NormalTok{, }\StringTok{"sentence\_id"}\NormalTok{))}
\end{Highlighting}
\end{Shaded}

The following code filters out invalid entity values, identifies unique
\texttt{doc\_id} and \texttt{entity} pairs, groups by brand and entity,
counts the frequency of each entity for each brand, sorts by brand by the
ascending number of appearances, and then selects the top 10 most
frequent entities per brand. A plot visualizes the frequency difference of the entities for each brand. For the meaning of the entities, check the
beginning of this chapter.

\begin{Shaded}
\begin{Highlighting}[]
\FunctionTok{library}\NormalTok{(ggplot2)}

\CommentTok{\# Count top 10 unique entities per brand}
\NormalTok{entity\_counts }\OtherTok{\textless{}{-}}\NormalTok{ sentence\_data }\SpecialCharTok{\%\textgreater{}\%}
  \FunctionTok{filter}\NormalTok{(}\SpecialCharTok{!}\FunctionTok{is.na}\NormalTok{(entity) }\SpecialCharTok{\&}\NormalTok{ entity }\SpecialCharTok{!=} \StringTok{""}\NormalTok{) }\SpecialCharTok{\%\textgreater{}\%}
  \FunctionTok{distinct}\NormalTok{(doc\_id, entity, }\AttributeTok{.keep\_all =} \ConstantTok{TRUE}\NormalTok{) }\SpecialCharTok{\%\textgreater{}\%}
  \FunctionTok{group\_by}\NormalTok{(Brand, entity) }\SpecialCharTok{\%\textgreater{}\%}
  \FunctionTok{summarize}\NormalTok{(}\AttributeTok{count =} \FunctionTok{n}\NormalTok{()) }\SpecialCharTok{\%\textgreater{}\%}
  \FunctionTok{arrange}\NormalTok{(Brand, }\FunctionTok{desc}\NormalTok{(count)) }\SpecialCharTok{\%\textgreater{}\%}
  \FunctionTok{group\_by}\NormalTok{(Brand) }\SpecialCharTok{\%\textgreater{}\%}
  \FunctionTok{slice\_max}\NormalTok{(count, }\AttributeTok{n =} \DecValTok{10}\NormalTok{) }

\CommentTok{\# Plot top entities for each brand}
\FunctionTok{ggplot}\NormalTok{(entity\_counts, }\FunctionTok{aes}\NormalTok{(}\AttributeTok{x =} \FunctionTok{reorder}\NormalTok{(entity, }\SpecialCharTok{{-}}\NormalTok{count), }
                          \AttributeTok{y =}\NormalTok{ count, }\AttributeTok{fill =}\NormalTok{ Brand)) }\SpecialCharTok{+}
  \FunctionTok{geom\_bar}\NormalTok{(}\AttributeTok{stat =} \StringTok{"identity"}\NormalTok{, }\AttributeTok{position =} \StringTok{"dodge"}\NormalTok{) }\SpecialCharTok{+}
  \FunctionTok{labs}\NormalTok{(}\AttributeTok{title =} \StringTok{"Most Frequent Entities for Each Brand"}\NormalTok{,}
       \AttributeTok{x =} \StringTok{"Entity"}\NormalTok{,}
       \AttributeTok{y =} \StringTok{"Count"}\NormalTok{,}
       \AttributeTok{fill =} \StringTok{"Brand"}\NormalTok{) }\SpecialCharTok{+}
  \FunctionTok{theme\_minimal}\NormalTok{() }\SpecialCharTok{+}
  \FunctionTok{coord\_flip}\NormalTok{()}
\end{Highlighting}
\end{Shaded}

\includegraphics{9.Chapter_new_files/figure-pdf/ner4-1.pdf}

Now, we want to to find which ORG or PRODUCT entities are present in
the negative reviews and how often. First, we filter for specific entities,
such as ``PRODUCT\_B'', ``ORG\_B'', ``ORG\_I'', then we capitalize the
tokens for both ORG and PRODUCT to ensure uniformity in the results. Then we
group the tokens by brand, count the occurrences, filter tokens with
more than 5 occurrences, and sort by descending counts. In the plot of
the most frequent entities for each of the two brands we can see how
often competitors are mentioned in these reviews. For example, in the
Samsung or Galaxy or Android reviews, there is no mention of an iPad or Apple product, while for Apple, Samsung is mentioned.
Amazon is mentioned for both brands because it is the vendor.

\begin{Shaded}
\begin{Highlighting}[]
\CommentTok{\# Filter and capitalize specific entities}
\NormalTok{op\_entities }\OtherTok{\textless{}{-}}\NormalTok{ sentence\_data }\SpecialCharTok{\%\textgreater{}\%}
  \FunctionTok{filter}\NormalTok{(entity }\SpecialCharTok{\%in\%} \FunctionTok{c}\NormalTok{(}\StringTok{"PRODUCT\_B"}\NormalTok{, }\StringTok{"ORG\_B"}\NormalTok{, }\StringTok{"ORG\_I"}\NormalTok{)) }\SpecialCharTok{\%\textgreater{}\%}
  \FunctionTok{mutate}\NormalTok{(}\AttributeTok{token =} \FunctionTok{ifelse}\NormalTok{(}\FunctionTok{startsWith}\NormalTok{(entity, }\StringTok{"ORG"}\NormalTok{) }\SpecialCharTok{|} 
    \FunctionTok{startsWith}\NormalTok{(entity, }\StringTok{"PRODUCT"}\NormalTok{), }\FunctionTok{toupper}\NormalTok{(token), token))}

\CommentTok{\# Count occurrences, filter, and sort}
\NormalTok{op\_counts }\OtherTok{\textless{}{-}}\NormalTok{ op\_entities }\SpecialCharTok{\%\textgreater{}\%}
  \FunctionTok{group\_by}\NormalTok{(token, Brand) }\SpecialCharTok{\%\textgreater{}\%}
  \FunctionTok{summarize}\NormalTok{(}\AttributeTok{count =} \FunctionTok{n}\NormalTok{()) }\SpecialCharTok{\%\textgreater{}\%}
  \FunctionTok{filter}\NormalTok{(count }\SpecialCharTok{\textgreater{}} \DecValTok{5}\NormalTok{) }\SpecialCharTok{\%\textgreater{}\%}
  \FunctionTok{arrange}\NormalTok{(}\FunctionTok{desc}\NormalTok{(count))}

\CommentTok{\# Plot most frequent entities by brand}
\FunctionTok{ggplot}\NormalTok{(op\_counts, }\FunctionTok{aes}\NormalTok{(}\AttributeTok{x =} \FunctionTok{reorder}\NormalTok{(token, }\SpecialCharTok{{-}}\NormalTok{count), }\AttributeTok{y =}\NormalTok{ count, }\AttributeTok{fill =}\NormalTok{ Brand)) }\SpecialCharTok{+}
  \FunctionTok{geom\_bar}\NormalTok{(}\AttributeTok{stat =} \StringTok{"identity"}\NormalTok{, }\AttributeTok{position =} \StringTok{"dodge"}\NormalTok{) }\SpecialCharTok{+}
  \FunctionTok{labs}\NormalTok{(}\AttributeTok{title =} \StringTok{"Most Frequent ORG/PROD Entities Subdivided by Brand"}\NormalTok{,}
       \AttributeTok{x =} \StringTok{"Organization"}\NormalTok{,}
       \AttributeTok{y =} \StringTok{"Count"}\NormalTok{,}
       \AttributeTok{fill =} \StringTok{"Brand"}\NormalTok{) }\SpecialCharTok{+}
  \FunctionTok{theme\_minimal}\NormalTok{() }\SpecialCharTok{+}
  \FunctionTok{coord\_flip}\NormalTok{()}
\end{Highlighting}
\end{Shaded}

\includegraphics{9.Chapter_new_files/figure-pdf/ner5-1.pdf}

\subsection{Extracting adjectives}\label{extracting-adjectives}

The \texttt{pos} field identifies which category a word, or its lemma,
belongs to, and uses a tag set that is generally well understood and
useful across various applications. For a more detailed treatment,
especially in English, \texttt{spaCy} uses the OntoNotes 5 version
of the Penn Treebank tag set, while other languages may use different
Treebanks. \footnote{spaCy's POS Tagging Documentation:
  https://spacy.io/docs/usage/pos-tagging\#pos-tagging-english} More
details can be found in spaCy's documentation. With this understanding
of the text structure, we can now proceed to extract specific
elements from the analyzed content.

Emotions and sensory perceptions are often expressed in language through adjectives, and it
is useful to understand this because they help us to
differentiate a brand and to detect underlying pain points. Repeated use
of adjectives in reviews, or their appearance in slogans (Ke and Wang
2013), might be persuasive (Ebaid 2018) for how a product or a
service is perceived and emphasize strengths or weaknesses. Adjectives are commonly
used in commerce, whether they appear in a simple basic form, or as comparatives
or superlatives (Hakkarainen 2016). They can elicit emotions and help people remember a specific attribute of a brand (Berger, Kim, and
Meyer 2021). We have selected the ADJ part of the reviews, grouped them by adjective
token and brand, counted the occurrences, filtered adjectives with more than
one occurrence, and then sorted by descending counts. After this, we have listed the
20 most frequent adjectives and the brands they are related to.

\begin{Shaded}
\begin{Highlighting}[]
\FunctionTok{library}\NormalTok{(tidyr)}

\CommentTok{\# Filter adjectives from sentence data and lowercase token}
\NormalTok{adj\_entities }\OtherTok{\textless{}{-}}\NormalTok{ sentence\_data }\SpecialCharTok{\%\textgreater{}\%}
  \FunctionTok{filter}\NormalTok{(pos }\SpecialCharTok{==} \StringTok{"ADJ"}\NormalTok{) }\SpecialCharTok{\%\textgreater{}\%}
  \FunctionTok{mutate}\NormalTok{(}\AttributeTok{token =} \FunctionTok{tolower}\NormalTok{(token))}

\CommentTok{\# Count occurrences of adjectives by brand}
\NormalTok{adj\_counts }\OtherTok{\textless{}{-}}\NormalTok{ adj\_entities }\SpecialCharTok{\%\textgreater{}\%}
  \FunctionTok{group\_by}\NormalTok{(token, Brand) }\SpecialCharTok{\%\textgreater{}\%}  
  \FunctionTok{summarize}\NormalTok{(}\AttributeTok{count =} \FunctionTok{n}\NormalTok{(), }\AttributeTok{.groups =} \StringTok{"drop"}\NormalTok{)  }\CommentTok{\# Count occurrences per brand}

\CommentTok{\# Create columns for Apple and Samsung brands and counts}
\NormalTok{adj\_counts\_wide }\OtherTok{\textless{}{-}}\NormalTok{ adj\_counts }\SpecialCharTok{\%\textgreater{}\%}
  \FunctionTok{pivot\_wider}\NormalTok{(}\AttributeTok{names\_from =}\NormalTok{ Brand, }\AttributeTok{values\_from =}\NormalTok{ count, }\AttributeTok{values\_fill =} \FunctionTok{list}\NormalTok{(}\AttributeTok{count =} \DecValTok{0}\NormalTok{)) }\SpecialCharTok{\%\textgreater{}\%}
  \FunctionTok{rename}\NormalTok{(}\StringTok{\textasciigrave{}}\AttributeTok{Brand Apple}\StringTok{\textasciigrave{}} \OtherTok{=}\NormalTok{ Apple, }\StringTok{\textasciigrave{}}\AttributeTok{Count Apple}\StringTok{\textasciigrave{}} \OtherTok{=}\NormalTok{ Apple, }
         \StringTok{\textasciigrave{}}\AttributeTok{Brand Samsung}\StringTok{\textasciigrave{}} \OtherTok{=}\NormalTok{ Samsung, }\StringTok{\textasciigrave{}}\AttributeTok{Count Samsung}\StringTok{\textasciigrave{}} \OtherTok{=}\NormalTok{ Samsung) }\SpecialCharTok{\%\textgreater{}\%}
  \FunctionTok{arrange}\NormalTok{(}\FunctionTok{desc}\NormalTok{(}\StringTok{\textasciigrave{}}\AttributeTok{Count Apple}\StringTok{\textasciigrave{}} \SpecialCharTok{+} \StringTok{\textasciigrave{}}\AttributeTok{Count Samsung}\StringTok{\textasciigrave{}}\NormalTok{))}

\CommentTok{\# Display the top 20 most frequent adjectives}
\NormalTok{top\_20\_adj }\OtherTok{\textless{}{-}}\NormalTok{ adj\_counts\_wide }\SpecialCharTok{\%\textgreater{}\%}
  \FunctionTok{slice\_max}\NormalTok{(}\AttributeTok{order\_by =}\NormalTok{ (}\StringTok{\textasciigrave{}}\AttributeTok{Count Apple}\StringTok{\textasciigrave{}} \SpecialCharTok{+} \StringTok{\textasciigrave{}}\AttributeTok{Count Samsung}\StringTok{\textasciigrave{}}\NormalTok{), }\AttributeTok{n =} \DecValTok{20}\NormalTok{)}

\CommentTok{\# Print results}
\FunctionTok{print}\NormalTok{(top\_20\_adj)}
\end{Highlighting}
\end{Shaded}

\begin{verbatim}
# A tibble: 20 x 3
   token        `Count Apple` `Count Samsung`
   <chr>                <int>           <int>
 1 other                   17              14
 2 new                     23               7
 3 same                     7              17
 4 bad                      4              16
 5 disappointed            11               8
 6 good                     8              11
 7 few                      6              12
 8 first                   11               6
 9 great                    6              11
10 more                    10               7
11 able                     4              10
12 many                     7               6
13 annoying                 1               9
14 long                     8               2
15 much                     2               8
16 sure                     5               4
17 unable                   6               3
18 dead                     5               3
19 most                     2               6
20 slow                     1               7
\end{verbatim}

Basic adjectives describe a quality in a simple form without making any
comparison. Comparative adjectives often end in ``-er'' and might appear
together with ``more'', ``much'', or ``less''. One such example is
``better'' or ``much better''. Superlative adjectives go to the extreme,
they often end in ``-est'' and are used with words like ``most'',
``least'', such in the ``best'' or ``least useful''. In our data set, of the
20 adjectives, two are comparative, ``more'' and ``much'', the
others are basic. As an exercise, we have left you the task of viewing all the
adjectives and categorizing them.

\subsection{Dependencies and POS
highlighting}\label{dependencies-and-pos-highlighting}

It is interesting and useful to obtain additional information from the
text and better understand how parts of speech contribute to meaning.
\texttt{spacyr}, in conjunction with the package \texttt{rsyntax}, help
us to plot a tree of POS, highlighting the dependencies between words. If you
want another example, you can substitute ``text1'' and assign a document of your choice to
``doc\_id''; see Figure~\ref{fig-ner7} to
visualize a syntactic-dependency parse tree and understand the
grammatical structure of a sentence. The words are related to one another
and this is indicated by the branches. The arrows points from a head
word to a dependent word. The numbers correspond to the word order in a
sentence to keep track of the order of appearance. The ROOT of the tree
is the main verb, representing the central action, while other words are
linked through different dependency relations. There are several labels
in the image, such as ``nsubjh'', ``dobj'', and ``punct'', which come from
the universal-dependencies (UD) annotation scheme. For example, ``nobj''
represents a nominal subject (who or what performs an action), ``dobj''
represents a direct object (the person or object that is the recipient of the action of a verb),
``punct'' is a punctuation mark. For a more complete list and definition of
the labels, see De Marneffe et al. (2021) . This type of analysis is
also suitable for chatbots, search engines, or grammar-correction tools.

\begin{Shaded}
\begin{Highlighting}[]
\FunctionTok{library}\NormalTok{(rsyntax)}

\CommentTok{\# Parse the text considering the dependencies}
\NormalTok{parsedtxt }\OtherTok{\textless{}{-}} \FunctionTok{spacy\_parse}\NormalTok{(data}\SpecialCharTok{$}\NormalTok{target, }\AttributeTok{dependency =} \ConstantTok{TRUE}\NormalTok{)}

\CommentTok{\# Plot the dependency tree for a specific document}
\FunctionTok{plot\_tree}\NormalTok{(}
  \FunctionTok{as\_tokenindex}\NormalTok{(parsedtxt),}
  \AttributeTok{doc\_id =} \StringTok{"text1"}\NormalTok{, }
\NormalTok{  token,}
\NormalTok{  lemma,}
\NormalTok{  pos,}
  \AttributeTok{viewer\_size =} \FunctionTok{c}\NormalTok{(}\DecValTok{110}\NormalTok{, }\DecValTok{150}\NormalTok{),}
  \AttributeTok{textsize =} \DecValTok{1}\NormalTok{,}
  \AttributeTok{spacing =} \DecValTok{2}\NormalTok{,}
  \AttributeTok{viewer\_mode =} \ConstantTok{FALSE}
\NormalTok{)}
\end{Highlighting}
\end{Shaded}

\begin{figure}[H]

\centering{

\includegraphics[width=1\textwidth,height=\textheight]{9.Chapter_new_files/figure-pdf/fig-ner7-1.pdf}

}

\caption{\label{fig-ner7}POS dependencies}

\end{figure}%

It is a useful to visualize some of the POS elements, especially those you are
interested in. In the next example, we will select the first text in the
data set and highlight the adjectives and nouns that refer to a number of
specific entities. Although printed using shades of gray, we have encoded the entities in blue, the
adjectives in red, and the nouns in green. Nouns are excellent
candidates for product features (Kauffmann et al. 2019) and will allow
you to capture the link between the adjectives and the
nouns immediately. For the entities we are trying to understand more about,
see the color-annotated text after the code.

\begin{Shaded}
\begin{Highlighting}[]
\FunctionTok{library}\NormalTok{(dplyr)}
\FunctionTok{library}\NormalTok{(htmltools)}

\CommentTok{\# Filter for a specific document}
\NormalTok{doc\_id\_nr }\OtherTok{\textless{}{-}} \StringTok{"text1"}
\NormalTok{parsed\_data }\OtherTok{\textless{}{-}}\NormalTok{ parsedtxt }\SpecialCharTok{\%\textgreater{}\%} \FunctionTok{filter}\NormalTok{(doc\_id }\SpecialCharTok{==}\NormalTok{ doc\_id\_nr)}

\CommentTok{\# Create a function to highlight annotations in text}
\NormalTok{highlight\_annotations }\OtherTok{\textless{}{-}} \ControlFlowTok{function}\NormalTok{(text, parsed\_data) \{}
\NormalTok{  tokens }\OtherTok{\textless{}{-}}\NormalTok{ parsed\_data}\SpecialCharTok{$}\NormalTok{token}
\NormalTok{  entities }\OtherTok{\textless{}{-}}\NormalTok{ parsed\_data}\SpecialCharTok{$}\NormalTok{entity}
\NormalTok{  pos\_tags }\OtherTok{\textless{}{-}}\NormalTok{ parsed\_data}\SpecialCharTok{$}\NormalTok{pos}
  
  \CommentTok{\# Initialize an empty string to store the highlighted text}
\NormalTok{  highlighted\_text }\OtherTok{\textless{}{-}} \StringTok{""}
  
  \ControlFlowTok{for}\NormalTok{ (i }\ControlFlowTok{in} \FunctionTok{seq\_along}\NormalTok{(tokens)) \{}
\NormalTok{    token }\OtherTok{\textless{}{-}}\NormalTok{ tokens[i]}
\NormalTok{    entity }\OtherTok{\textless{}{-}}\NormalTok{ entities[i]}
\NormalTok{    pos\_tag }\OtherTok{\textless{}{-}}\NormalTok{ pos\_tags[i]}
    
    \CommentTok{\# Highlight based on entity and POS tag}
    \ControlFlowTok{if}\NormalTok{ (entity }\SpecialCharTok{!=} \StringTok{""}\NormalTok{) \{}
\NormalTok{      highlighted\_text }\OtherTok{\textless{}{-}} \FunctionTok{paste0}\NormalTok{(highlighted\_text, }
                      \StringTok{"\textless{}span style=\textquotesingle{}color: blue;\textquotesingle{}\textgreater{}"}\NormalTok{, token, }\StringTok{"\textless{}/span\textgreater{}"}\NormalTok{)}
\NormalTok{    \} }\ControlFlowTok{else} \ControlFlowTok{if}\NormalTok{ (pos\_tag }\SpecialCharTok{==} \StringTok{"ADJ"}\NormalTok{) \{}
\NormalTok{      highlighted\_text }\OtherTok{\textless{}{-}} \FunctionTok{paste0}\NormalTok{(highlighted\_text, }
                      \StringTok{"\textless{}span style=\textquotesingle{}color: red;\textquotesingle{}\textgreater{}"}\NormalTok{, token, }\StringTok{"\textless{}/span\textgreater{}"}\NormalTok{)}
\NormalTok{    \} }\ControlFlowTok{else} \ControlFlowTok{if}\NormalTok{ (pos\_tag }\SpecialCharTok{==} \StringTok{"NOUN"}\NormalTok{) \{}
\NormalTok{      highlighted\_text }\OtherTok{\textless{}{-}} \FunctionTok{paste0}\NormalTok{(highlighted\_text, }
                      \StringTok{"\textless{}span style=\textquotesingle{}color: green;\textquotesingle{}\textgreater{}"}\NormalTok{, token, }\StringTok{"\textless{}/span\textgreater{}"}\NormalTok{)}
\NormalTok{    \} }\ControlFlowTok{else}\NormalTok{ \{}
\NormalTok{      highlighted\_text }\OtherTok{\textless{}{-}} \FunctionTok{paste0}\NormalTok{(highlighted\_text, token)}
\NormalTok{    \}}
    
    \CommentTok{\# Add a space after each token}
\NormalTok{    highlighted\_text }\OtherTok{\textless{}{-}} \FunctionTok{paste0}\NormalTok{(highlighted\_text, }\StringTok{" "}\NormalTok{)}
\NormalTok{  \}}

  \FunctionTok{return}\NormalTok{(highlighted\_text)}
\NormalTok{\}}

\CommentTok{\# Apply the function to the parsed text}
\NormalTok{highlighted\_text }\OtherTok{\textless{}{-}} \FunctionTok{highlight\_annotations}\NormalTok{(text, parsed\_data)}

\CommentTok{\# Display inline}
\FunctionTok{browsable}\NormalTok{(}\FunctionTok{HTML}\NormalTok{(highlighted\_text))}
\end{Highlighting}
\end{Shaded}

{\color{darkgreen} Display} of {\color{darkgreen} brand} {\color{red} new} {\color{darkgreen} unit} bought from {\color{blue} amazon.com} has an {\color{red} oily like} {\color{darkgreen} film spot} under the {\color{darkgreen} screen}. It seems to be a {\color{darkgreen} manufacturing defect}, according to {\color{darkgreen} internet research} {\color{darkgreen} lots} of {\color{darkgreen} users} do experience the {\color{red} same} {\color{darkgreen} problem} with their {\color{darkgreen} units}. Waiting for {\color{blue} Samsung} to fix that {\color{darkgreen} problem} or {\color{darkgreen} recall} of the {\color{red} affected} {\color{darkgreen} units}, cause {\color{red} most} {\color{darkgreen} customers} still have this {\color{darkgreen} screen problem}, even after they have had their {\color{darkgreen} units} exchanged. {\color{darkgreen} Btw} that {\color{darkgreen} problem} on my {\color{darkgreen} tab} is {\color{red} visible} even when the {\color{darkgreen} unit} is {\color{red} cold} - so I do not think that it is necessarily {\color{darkgreen} heat} related. Otherwise, it is really a {\color{red} great} {\color{blue} Tab} - but {\color{red} sorry} {\color{blue} Samsung} - for {\color{blue} 600 \$} I would like to own a {\color{darkgreen} unit} without a {\color{darkgreen} flaw}. {\color{darkgreen} Hope} that there will be soon a {\color{darkgreen} solution} or a {\color{red} worldwide generous} {\color{darkgreen} recall / replacement offer} to satisfy {\color{blue} Galaxy users \& fans}.

Explore this review and run the code again to visualize another review
by simply substituting ``text1'' with another value, such as ``text2'',
``text10'', or ``tex42''. You can create several combinations by
extracting not only the adjectives and nouns, but also adverbs. Adverbs can
be classified as being of time, frequency, location, manner, degree, or
as conjunctive adverbs. In sentiment analysis, the adverbs of degree, such
as ``extremely'', ``absolutely'', ``really'', and so on, are the most
relevant ones. These adverbs have a further classification, as found in
(Benamara et al. 2007), and a few examples of each class are given below:

\begin{itemize}
\tightlist
\item
  Adverbs of affirmation: certainly, definitely, absolutely,
  surely, indeed, undoubtedly, clearly, obviously
\item
  Adverbs of doubt,: possibly, maybe, possibly, apparently,
  likely, conceivably, arguably, supposedly, uncertainly
\item
  Strong intensifying adverbs: extremely, incredibly,
  absolutely, totally, completely, remarkably, exceptionally, amazingly,
  terribly, deeply
\item
  Weak intensifying adverbs: somewhat, fairly, slightly,
  moderately, rather, mildly, partly, relatively
\item
  Negation and minimizers are those that negate or diminish sentiments: never, nowhere, nothing, hardly, barely, rarely, little,
  infrequently
\end{itemize}

Adverbs, ADV, can be extracted in a similar manner to adjectives, and
the detection of both adjectives and adverbs in text (Roy et al. 2015)
can be used to shape targeted marketing messages.

\subsection{Tracking adjective usage over
time}\label{tracking-adjective-usage-over-time}

Marketers are interested in assessing the performance of their
efforts by analyzing the prevalence of adjectives in consumer reviews
before, during, and after marketing campaigns or product launches. A
successful campaign might be reflected in an increase in the use of positive
adjectives and a decrease in negative ones. We next demonstrate the
tracking over time of the frequency of negative reviews for a brand, together with
those containing a competitor. If the information is
extracted properly from the dataset, such tracking can also be helpful in
acquiring seasonal insights and supporting crisis management. The
following code filters the \texttt{comparison} dataset to include only
adjectives related to two brands (Apple and Samsung). Then it groups and
counts the adjectives by date and brand, creates a line chart to display
the adjective usage over time, and adds a smoothing curve to show the trends.

\begin{Shaded}
\begin{Highlighting}[]
\CommentTok{\# Filter for adjectives and the two brands}
\NormalTok{adjective\_data }\OtherTok{\textless{}{-}}\NormalTok{ sentence\_data }\SpecialCharTok{\%\textgreater{}\%}
  \FunctionTok{filter}\NormalTok{(pos }\SpecialCharTok{==} \StringTok{"ADJ"} \SpecialCharTok{\&}\NormalTok{ Brand }\SpecialCharTok{\%in\%} \FunctionTok{c}\NormalTok{(}\StringTok{"Apple"}\NormalTok{, }\StringTok{"Samsung"}\NormalTok{))}

\CommentTok{\# Group by date and brand, then count the number of adjectives}
\NormalTok{adjective\_counts }\OtherTok{\textless{}{-}}\NormalTok{ adjective\_data }\SpecialCharTok{\%\textgreater{}\%}
  \FunctionTok{group\_by}\NormalTok{(R\_date, Brand) }\SpecialCharTok{\%\textgreater{}\%}
  \FunctionTok{summarize}\NormalTok{(}\AttributeTok{count =} \FunctionTok{n}\NormalTok{()) }\SpecialCharTok{\%\textgreater{}\%}
  \FunctionTok{ungroup}\NormalTok{()}

\CommentTok{\# Plotting the data with both lines and smoothing}
\FunctionTok{ggplot}\NormalTok{(adjective\_counts, }\FunctionTok{aes}\NormalTok{(}\AttributeTok{x =}\NormalTok{ R\_date, }\AttributeTok{y =}\NormalTok{ count, }\AttributeTok{color =}\NormalTok{ Brand, }\AttributeTok{linetype =}\NormalTok{ Brand)) }\SpecialCharTok{+}
  \FunctionTok{geom\_line}\NormalTok{() }\SpecialCharTok{+}
  \FunctionTok{geom\_smooth}\NormalTok{(}\AttributeTok{method =} \StringTok{"loess"}\NormalTok{, }\AttributeTok{span =} \FloatTok{0.3}\NormalTok{, }\AttributeTok{se =} \ConstantTok{FALSE}\NormalTok{) }\SpecialCharTok{+}
  \FunctionTok{labs}\NormalTok{(}\AttributeTok{title =} \StringTok{"Number of Adjectives Over Time for Apple and Samsung"}\NormalTok{,}
       \AttributeTok{x =} \StringTok{"Date"}\NormalTok{,}
       \AttributeTok{y =} \StringTok{"Number of Adjectives"}\NormalTok{) }\SpecialCharTok{+}
  \FunctionTok{theme\_minimal}\NormalTok{() }\SpecialCharTok{+}
  \FunctionTok{theme}\NormalTok{(}\AttributeTok{axis.text.x =} \FunctionTok{element\_text}\NormalTok{(}\AttributeTok{angle =} \DecValTok{45}\NormalTok{, }\AttributeTok{hjust =} \DecValTok{1}\NormalTok{)) }\SpecialCharTok{+}
  \FunctionTok{guides}\NormalTok{(}\AttributeTok{color =} \FunctionTok{guide\_legend}\NormalTok{(}\AttributeTok{override.aes =} \FunctionTok{list}\NormalTok{(}\AttributeTok{size =} \DecValTok{1}\NormalTok{)))}
\end{Highlighting}
\end{Shaded}

\begin{figure}[H]

\centering{

\includegraphics{9.Chapter_new_files/figure-pdf/fig-ner9-1.pdf}

}

\caption{\label{fig-ner9}}

\end{figure}%

The graph in Figure~\ref{fig-ner9} plots the frequency of adjectives
associated with two brands, Apple and Samsung. The trend for Apple shows
fluctuations, with a general increase in the use of adjectives from 2010
to around mid-2012. A smooth trend line can also be seen. Samsung
showed more variability between mid-2011 and late 2012. There is a
general increase in variability over time, with peaks becoming more
extreme. This might indicate more varied public and media reactions
through a spike in adjective usage. Consumers expressed negative feelings
about Apple's products through the steady use of adjectives. Given that
customers express their anger differently for the two brands, marketers
can use this information to shape their communication strategies. The language used
can be targeted at different segments of the population. It is possible to modify the
previous code to add more brands, use a different smoothing method (for
example a linear model, ``lm'') or use a stacked bar chart instead:
replace \texttt{geom\_line()} with
\texttt{geom\_bar(stat\ =\ "identity",\ position\ =\ "stack")}.

\subsection{Language complexity}\label{language-complexity}

Sometimes researchers are interested in studying the syntactic
complexity in relation to the language proficiency shown
in texts, and the overall (or ``objective'') proficiency in a target
language (Ortega 2003). When conducting these measurements, it is important to
compare them against a benchmark or another unit, for example, unit A
compared to unit B on a specific index can provide insights into
accuracy and applicability (Szmrecsanyi 2004). There are several metrics
used to calculate syntactic complexity, and they are mainly used in
linguistics, language acquisition, and computational linguistics. In
this section, we demonstrate the use of a simple metric, which is
based on the measurement of the density of verbs called Fichtner's C
(Fichtner 1980; Samples, Ireland, and Kraczon 2024), which is defined
as:

\[
\text{Fichtner's C} = \frac{\text{Number of Verbs}}{\text{Number of Sentences}} \times \frac{\text{Number of Words}}{\text{Number of Sentences}}
\]

Using syntactic complexity in marketing allows for a detailed comparison of
the language used to describe different brands. By analyzing metrics
like Fichtner's C, marketers can gain insights into how brands are
perceived and what level of complexity is used to describe them, thus
tailoring their communication styles to specific audiences. The code
below computes the index for Fichtner's C and compares the linguistic complexity
between the brands on a table.

\begin{Shaded}
\begin{Highlighting}[]
\CommentTok{\# Load necessary libraries}
\FunctionTok{library}\NormalTok{(dplyr)}
\FunctionTok{library}\NormalTok{(ggplot2)}

\CommentTok{\# Calculate linguistic complexity for each document and brand}
\NormalTok{syn\_com }\OtherTok{\textless{}{-}}\NormalTok{ sentence\_data }\SpecialCharTok{\%\textgreater{}\%}
  \FunctionTok{group\_by}\NormalTok{(doc\_id, Brand) }\SpecialCharTok{\%\textgreater{}\%}
  \FunctionTok{summarize}\NormalTok{(}\AttributeTok{verb\_num =} \FunctionTok{sum}\NormalTok{(pos }\SpecialCharTok{==} \StringTok{"VERB"}\NormalTok{),}
            \AttributeTok{sent\_num =} \FunctionTok{n\_distinct}\NormalTok{(sentence\_id),}
            \AttributeTok{word\_num =} \FunctionTok{n}\NormalTok{()) }\SpecialCharTok{\%\textgreater{}\%}
  \FunctionTok{mutate}\NormalTok{(}\AttributeTok{F\_C =}\NormalTok{ (verb\_num }\SpecialCharTok{/}\NormalTok{ sent\_num) }\SpecialCharTok{*}\NormalTok{ (word\_num }\SpecialCharTok{/}\NormalTok{ sent\_num)) }\SpecialCharTok{\%\textgreater{}\%}
  \FunctionTok{ungroup}\NormalTok{()}

\CommentTok{\# Compare linguistic complexity between the two brands}
\NormalTok{brand\_comparison }\OtherTok{\textless{}{-}}\NormalTok{ syn\_com }\SpecialCharTok{\%\textgreater{}\%}
  \FunctionTok{group\_by}\NormalTok{(Brand) }\SpecialCharTok{\%\textgreater{}\%}
  \FunctionTok{summarize}\NormalTok{(}\AttributeTok{avg\_F\_C =} \FunctionTok{mean}\NormalTok{(F\_C, }\AttributeTok{na.rm =} \ConstantTok{TRUE}\NormalTok{),}
            \AttributeTok{sd\_F\_C =} \FunctionTok{sd}\NormalTok{(F\_C, }\AttributeTok{na.rm =} \ConstantTok{TRUE}\NormalTok{),}
            \AttributeTok{median\_F\_C =} \FunctionTok{median}\NormalTok{(F\_C, }\AttributeTok{na.rm =} \ConstantTok{TRUE}\NormalTok{),}
            \AttributeTok{n =} \FunctionTok{n}\NormalTok{())}

\CommentTok{\# Display the result}
\FunctionTok{print}\NormalTok{(brand\_comparison)}
\end{Highlighting}
\end{Shaded}

\begin{verbatim}
# A tibble: 2 x 5
  Brand   avg_F_C sd_F_C median_F_C     n
  <chr>     <dbl>  <dbl>      <dbl> <int>
1 Apple      55.2   89.0       37.7    87
2 Samsung    42.8   44.1       26.5    71
\end{verbatim}

From the table, we can see that reviews containing Apple
exhibit, on average, higher linguistic complexity, with a greater
than average Fichtner's C index (\texttt{avg\_F\_C}) and a greater
variability than Samsung, as shown by the higher standard deviation
(\texttt{sd\_F\_C}). This could imply that content about Apple ranges
from very simple to very complex, possibly depending on the type of
document or context. In contrast, documents about Samsung are more
consistent in terms of linguistic complexity, as indicated by the lower
standard deviation (\texttt{sd\_F\_C}). This suggests that reviews about
Samsung maintain a more uniform style across its documents. Both brands
have median values (\texttt{median\_F\_C}) values lower than their
respective averages (\texttt{avg\_F\_C}), indicating a right-skewed
distribution. This means that most reviews have moderate levels of complexity but
there are some documents with very high complexity which raise the
average. Notice that differences in reviews (in this case across brands)
can also reflect the heterogeneity in language used by the
clientele and not necessarily is attributable to the brand as such. To
visualize this, we can also plot the linguistic complexity measures.

\begin{Shaded}
\begin{Highlighting}[]
\CommentTok{\# Plot the density of linguistic complexity for the two brands}
\FunctionTok{ggplot}\NormalTok{(syn\_com, }\FunctionTok{aes}\NormalTok{(}\AttributeTok{x =}\NormalTok{ F\_C, }\AttributeTok{fill =}\NormalTok{ Brand)) }\SpecialCharTok{+}
  \FunctionTok{geom\_density}\NormalTok{(}\AttributeTok{alpha =} \FloatTok{0.5}\NormalTok{) }\SpecialCharTok{+}
  \FunctionTok{labs}\NormalTok{(}\AttributeTok{title =} \StringTok{"Comparison of Linguistic Complexity for Apple and Samsung"}\NormalTok{,}
       \AttributeTok{x =} \StringTok{"Linguistic Complexity (F\_C)"}\NormalTok{,}
       \AttributeTok{y =} \StringTok{"Density"}\NormalTok{) }\SpecialCharTok{+}
  \FunctionTok{theme\_minimal}\NormalTok{() }\SpecialCharTok{+}
  \FunctionTok{scale\_fill\_manual}\NormalTok{(}\AttributeTok{values =} \FunctionTok{c}\NormalTok{(}\StringTok{"Apple"} \OtherTok{=} \StringTok{"blue"}\NormalTok{, }\StringTok{"Samsung"} \OtherTok{=} \StringTok{"red"}\NormalTok{)) }\SpecialCharTok{+}
  \FunctionTok{theme}\NormalTok{(}\AttributeTok{legend.title =} \FunctionTok{element\_blank}\NormalTok{())}
\end{Highlighting}
\end{Shaded}

\includegraphics{9.Chapter_new_files/figure-pdf/ner11-1.pdf}

The visualization confirms out previous discussion. At the lower end of
the linguistic complexity scale, both brands show a peak, and this shows
that a larger proportion of their documents have low linguistic complexity.
Apple reviews have a higher peak density compared to Samsung, meaning
they have greater linguistic complexity. The brands have long tails which reach higher levels of
linguistic complexity, with Samsung having a more pronounced
tail. This indicates that some documents have higher levels of complexity. Across
the negative reviews, customers refer to Apple with more consistent
linguistic complexity, so its documents are more homogenous. Customers who
write about Samsung in a negative context use a slightly more varied
writing style. Nevertheless, both brands, when they are described in
negative terms, have a low level of linguistic complexity, indicating a precise
focus and clear, straightforward communication.

\subsection{A final technical note}\label{a-final-technical-note}

Several advanced techniques can be employed to enhance an NER task. An example is the use of consolidation functions to
merge multi-word entities into a single token, or the concatenation of
words into noun phrases, also called noun chunks. \texttt{SpaCy} can use
language models which are not limited to English, but also include
German, Spanish, Portugues, French, Italian, and Dutch, among others.
\texttt{spacyr} can be integrated with other R text-analysis packages,
for example, \texttt{quanteda}, or \texttt{tidytext}, thus simplifying
text manipulation and analysis tasks.

Additionally, as we have seen in Chapter~\ref{sec-textual-data}, regular
expressions can be helpful when extracting preposition-noun combinations
from the corpus, and checking the most frequent noun for each preposition
or other parts of speech. A clear understanding of the annotations
and the meanings of the tags at different levels, as well as the tag sets
present in various languages, is essential because these could influence
your approach to NER and text analysis in general.

Using Python in the background of an R session can be very
memory-intensive. Once your analysis is done, it is strongly recommended to
use the \texttt{spacy\_finalize()} function which terminates Python
and frees up the memory it was using.

\begin{Shaded}
\begin{Highlighting}[]
\CommentTok{\# Clean up and release spaCy resources}
\FunctionTok{spacy\_finalize}\NormalTok{()}
\end{Highlighting}
\end{Shaded}

\section{Summarizing}\label{summarization}

The word summary relates to something which is brief
and concise. The word comes from the Latin word \textit{summa}, which means
``totality, gist''. One will find summaries everywhere, such as a
summary of a book, in the workplace, in the minutes of meetings, in news
articles, and in customer reviews. Because we focus on the latter,
we have to understand that this allows for an aggregated view of opinions,
delivering something similar to''the wisdom of the crowd''. Many texts
are too long to read, and we often have too little time, and so the comment ``too long, did not read''. The abbreviation for this is TL;DR (again
a shorter version of something), which is commonly used in online
discussions, social media, etc., indicating content that is lengthy and that a
summary is provided for whoever does not have time to read the whole
text. TL;DR can be used as a label for the summary itself, or as a
comment to indicate that the text is too long. Summaries are also
useful for information retention and enhance the ability to learn
something (Reder and Anderson 1980), implying that they conveying
information in a very abridged form, creating a shortened version that
retains only the essential information and main points. But
how much can one summarize? In NLP, summarizing can, generally, be achieved by
using three approaches: extractive (Moratanch and Chitrakala 2017),
abstractive (Gupta and Gupta 2019), or hybrid (El-Kassas et al.
2021). The first approach uses a selection (and extraction) of key
sentences, or phrases, from the original text that act as a summary.
This selection is typically based on statistical features such as
importance scores, term frequency, or position. Abstractive summarizing is a
generative process that uses the rephrasing or paraphrasing of words which
might not exist in the original text. This last approach involves
understanding the context of a text and producing an abridged version that
captures the gist of the original. The hybrid method is a
combination of the first two. Below we will focus on extractive
summarization because it can be accomplished relatively quickly, although
it might be computationally intensive. It also fits nicely with the NER
elements we have discussed so far.

\subsection{Summarizing in R}\label{summarization-in-r}

There are a few packages that help with the task of summarizing in R,
either directly or indirectly. One is \texttt{lexRankr}, another is
LSAfun, and yet another one is \texttt{textrank}. The simplest one to
work with is \texttt{lexRankr}. The implementation of the algorithm in R
uses a stochastic graph-based method to assess the importance of a unit of text, such as a sentence, in a corpus (Erkan and Radev 2004). The basic
idea is that if the algorithm finds sentences with are similar to other
sentences, it means that they are important. In this graph-base
approach, sentences are nodes and edges represent the similarity between
sentences. To illustrate, let us take the ``data'' data frame and apply
the algorithm.

\begin{Shaded}
\begin{Highlighting}[]
\FunctionTok{library}\NormalTok{(lexRankr)}
\FunctionTok{library}\NormalTok{(dplyr)}

\CommentTok{\# Tokenize text data into sentences}
\NormalTok{tokenized\_data }\OtherTok{\textless{}{-}}\NormalTok{ data }\SpecialCharTok{\%\textgreater{}\%}
  \FunctionTok{unnest\_sentences}\NormalTok{(sentence, target)}

\CommentTok{\# Apply LexRank to rank sentences based on importance}
\NormalTok{ranked\_sentences }\OtherTok{\textless{}{-}} \FunctionTok{bind\_lexrank}\NormalTok{(tokenized\_data, sentence,}
\NormalTok{                                 doc\_id, }\AttributeTok{level =} \StringTok{"sentences"}\NormalTok{)}


\CommentTok{\# Filter and get top 3 sentences for Apple}
\NormalTok{top\_apple }\OtherTok{\textless{}{-}}\NormalTok{ ranked\_sentences }\SpecialCharTok{\%\textgreater{}\%}
  \FunctionTok{filter}\NormalTok{(Brand }\SpecialCharTok{==} \StringTok{"Apple"}\NormalTok{) }\SpecialCharTok{\%\textgreater{}\%}
  \FunctionTok{arrange}\NormalTok{(}\FunctionTok{desc}\NormalTok{(lexrank)) }\SpecialCharTok{\%\textgreater{}\%}
  \FunctionTok{slice}\NormalTok{(}\DecValTok{1}\SpecialCharTok{:}\DecValTok{3}\NormalTok{)}

\CommentTok{\# Filter and get top 3 sentences for Samsung}
\NormalTok{top\_samsung }\OtherTok{\textless{}{-}}\NormalTok{ ranked\_sentences }\SpecialCharTok{\%\textgreater{}\%}
  \FunctionTok{filter}\NormalTok{(Brand }\SpecialCharTok{==} \StringTok{"Samsung"}\NormalTok{) }\SpecialCharTok{\%\textgreater{}\%}
  \FunctionTok{arrange}\NormalTok{(}\FunctionTok{desc}\NormalTok{(lexrank)) }\SpecialCharTok{\%\textgreater{}\%}
  \FunctionTok{slice}\NormalTok{(}\DecValTok{1}\SpecialCharTok{:}\DecValTok{3}\NormalTok{)}

\CommentTok{\# Display the results, for Apple}
\FunctionTok{print}\NormalTok{(top\_apple}\SpecialCharTok{$}\NormalTok{sentence)}

\CommentTok{\# Display the results, for Samsung}
\FunctionTok{print}\NormalTok{(top\_samsung}\SpecialCharTok{$}\NormalTok{sentence)}
\end{Highlighting}
\end{Shaded}

\begin{verbatim}
[1] "the ipad in ?"                            
[2] "But it is not an iPad2."                  
[3] "No more apple products for us after this."
\end{verbatim}

\begin{verbatim}
[1] "Its slower than ipad."                                                                   
[2] "I did a lot of research on this tablet before buying and saw nothing about this problem."
[3] "I will never buy a samsung product again!"                                               
\end{verbatim}

Considering that the data set was selected from negative reviews, the
customers who talk about the Apple products express themselves in
negative terms but not in a vehement way. The customers who talk about
the Samsung tablets, however, seem to prefer Apple products.

To illustrate the use of another package to summarize text let us take a
look at \texttt{LSAfun}, which uses latent semantic analysis (Gong and Liu 2001). The basic idea is to decompose the text
into sentences and then choose a few \texttt{k}-candidate sentences to
form the summary. Let us imagine a simple example, where we extract two
sentences from the negative reviews data set, and try to understand the customers' reviews for
both Apple and Samsung products from two sentences each. The
function that performs this task is called \texttt{genericSummary()} and
has several arguments. The first, is the set of documents to analyze.
The second is the number of sentences it needs to extract. And the
third, called \texttt{split}, is a character vector which allows you to
chose what symbols indicate the end of a sentence. \texttt{min}
indicates the minimum amount of words a sentence should contain. Here we
will use only the first two arguments.

\begin{Shaded}
\begin{Highlighting}[]
\CommentTok{\#install.packages("LSAfun")}
\FunctionTok{library}\NormalTok{(LSAfun)}
\FunctionTok{library}\NormalTok{(dplyr)}

\CommentTok{\# We can recycle the sent\_df dataframe because it contains all the elements we need.}

\NormalTok{apple\_data }\OtherTok{\textless{}{-}}\NormalTok{ sent\_df }\SpecialCharTok{\%\textgreater{}\%} 
  \FunctionTok{filter}\NormalTok{(Brand }\SpecialCharTok{==} \StringTok{"Apple"}\NormalTok{)}

\NormalTok{samsung\_data }\OtherTok{\textless{}{-}}\NormalTok{ sent\_df }\SpecialCharTok{\%\textgreater{}\%} 
  \FunctionTok{filter}\NormalTok{(Brand }\SpecialCharTok{==} \StringTok{"Samsung"}\NormalTok{)}

\CommentTok{\# Summarize to 2 sentences}
\NormalTok{summary\_apple }\OtherTok{\textless{}{-}} \FunctionTok{genericSummary}\NormalTok{(apple\_data}\SpecialCharTok{$}\NormalTok{sentence, }\AttributeTok{k =} \DecValTok{2}\NormalTok{)  }
\NormalTok{summary\_samsung }\OtherTok{\textless{}{-}} \FunctionTok{genericSummary}\NormalTok{(samsung\_data}\SpecialCharTok{$}\NormalTok{sentence, }\AttributeTok{k =} \DecValTok{2}\NormalTok{)}

\CommentTok{\# Display the results, for Apple}
\FunctionTok{print}\NormalTok{(summary\_apple)}

\CommentTok{\# Display the results, for Samsung}
\FunctionTok{print}\NormalTok{(summary\_samsung)}
\end{Highlighting}
\end{Shaded}

\begin{verbatim}
[1] "I don't know why this P.O.S. has trouble with so many sites but here are some examples of things I've run into (and these are condensed): 1) just now while on Tripadvisor.com, I was trying to access the second page on the list of 'Local Attractions' and it froze while loading, showing the message \"Updating your results\" 2) I am unable to scroll down on certain drop-down boxes on websites 3) occasionally while doing a Google search I am unable to modify my search terms and it just defaults back to my original search no matter how I attempt to enter it"
[2] "when i looked at the box, the original serial number was scratched off the barcode label, telling me that the vendor had sent another (broken) device instead of the one that was originally in the box"                                                                                                                                                                                                                                                                                                                                                                       
\end{verbatim}

\begin{verbatim}
[1] "Im very disappointed of the touch screen sensitivity, nothing compared with the ipod touch i had, there are sometimes that i have to touch more than 4 or 5 times to get a response It is a big oroblem when i try to logouton facebook always is the same i have to zoom out the screen to make big buttons and not only in facebook, in all web pages is the same, i tried many web browsers and tje tablet is always calibrated but always i have problems with the screen i think this is a software problem, because there is a program to test the sensitivity and it seems that the screens response ok, I hope they will fix this soon because this is very annoying, I like other companies than apple make tablets because I was tired of itunes and all that, that is why i tried other brand."
[2] "Of course it shouldn't be hard to change this to my preferred language (English UK) but it was sort of annoying to see that and have to correct it."                                                                                                                                                                                                                                                                                                                                                                                                                                                                                                                                                                                                                                                      
\end{verbatim}

With this approach, we can see that the documents extracted are much
longer, but are they really representative of the two brands?

Another graph-based ranking method (Mihalcea and Tarau 2004) used to extract
sentences or keywords is \texttt{textrank} (Wijffels 2020). The nodes of
the graph represent a sentence, and the edges between them are
computations of the number of words adjacent nodes have in common. This
structure is passed on to the PageRank algorithm (Xing and Ghorbani
2004; implemented in Google's search engine, Brin and Page 1998, and
named after one of its founders, Larry Page) to find the most important
sentences which will form the summary of the text. In other words, the
summary is made out of the most important sentences. A summary can also be
created by using keywords. These represent documents and have the
capability to distill their content and express concepts, themes, or
topics in just a few words, making the process of information indexing,
searching, and retrieval easier. In the \texttt{textranking} algorithm, a
word network is built in a similar way to those that you have already encountered in
Chapter~\ref{sec-sentiment-analysis},
Section~\ref{sec-sentiment-word-networks}, for example. The link between
words becomes important when two words frequently follow one another.
In this case, PageRank is used as well, and the top third of
the keywords is considered important. Finding which words are
relevant can be done through the POS tagging we have seen in the
NER part of this chapter using \texttt{spacyr}. Let us set this up
and extract the important keywords. We have created a table that contains the
sentences, lemmas, and parts of speech. We can look at the newly created
data frame which contains all the information we need by selecting the fields ``sentence\_id'', ``lemma'', and ``upos''.

\begin{Shaded}
\begin{Highlighting}[]
\CommentTok{\# Add a unique identifier to each sentence and select relevant columns}
\NormalTok{sent\_df\_all }\OtherTok{\textless{}{-}}\NormalTok{ sentence\_data }\SpecialCharTok{\%\textgreater{}\%}
  \FunctionTok{mutate}\NormalTok{(}\AttributeTok{textrank\_id =} \FunctionTok{paste}\NormalTok{(doc\_id, sentence\_id, }\AttributeTok{sep =} \StringTok{"\_"}\NormalTok{)) }\SpecialCharTok{\%\textgreater{}\%}
  \FunctionTok{select}\NormalTok{(textrank\_id, doc\_id, sentence\_id, sentence, token, lemma, pos, Brand)}

\CommentTok{\# Filter the data for sentences related to Apple}
\NormalTok{sent\_df\_apple }\OtherTok{\textless{}{-}}\NormalTok{ sent\_df\_all }\SpecialCharTok{\%\textgreater{}\%}
  \FunctionTok{filter}\NormalTok{(Brand }\SpecialCharTok{==} \StringTok{"Apple"}\NormalTok{)}

\CommentTok{\# Filter the data for sentences related to Samsung}
\NormalTok{sent\_df\_smsng }\OtherTok{\textless{}{-}}\NormalTok{ sent\_df\_all }\SpecialCharTok{\%\textgreater{}\%}
  \FunctionTok{filter}\NormalTok{(Brand }\SpecialCharTok{==} \StringTok{"Samsung"}\NormalTok{)}

\CommentTok{\# Display the first 10 rows of selected columns for Samsung sentences}
\FunctionTok{head}\NormalTok{(sent\_df\_smsng[, }\FunctionTok{c}\NormalTok{(}\StringTok{"sentence\_id"}\NormalTok{, }\StringTok{"lemma"}\NormalTok{, }\StringTok{"pos"}\NormalTok{)], }\DecValTok{10}\NormalTok{)}
\end{Highlighting}
\end{Shaded}

\begin{verbatim}
   sentence_id      lemma   pos
1            1    display  NOUN
2            1         of   ADP
3            1      brand  NOUN
4            1        new   ADJ
5            1       unit  NOUN
6            1        buy  VERB
7            1       from   ADP
8            1 amazon.com PROPN
9            1       have  VERB
10           1         an   DET
\end{verbatim}

Now, we can extract the most important keywords for both brands. For this
example we will choose nouns and adjectives and subset the datasets for
multi-word frequencies. For space reasons, we have chosen different
frequencies for the brands, but it is possible to play around and select a
frequency threshold of \(>0\).

\begin{Shaded}
\begin{Highlighting}[]
\FunctionTok{library}\NormalTok{(textrank)  }\CommentTok{\# Load library}

\CommentTok{\# Extract keywords for Apple (nouns/adjectives)}
\NormalTok{tr\_kw\_a }\OtherTok{\textless{}{-}} \FunctionTok{textrank\_keywords}\NormalTok{(sent\_df\_apple}\SpecialCharTok{$}\NormalTok{lemma, }
                             \AttributeTok{relevant =}\NormalTok{ sent\_df\_apple}\SpecialCharTok{$}\NormalTok{pos }\SpecialCharTok{\%in\%} \FunctionTok{c}\NormalTok{(}\StringTok{"NOUN"}\NormalTok{, }\StringTok{"ADJ"}\NormalTok{))}
\CommentTok{\# Filter: multi{-}word, freq \textgreater{}= 2}
\FunctionTok{subset}\NormalTok{(tr\_kw\_a}\SpecialCharTok{$}\NormalTok{keywords, ngram }\SpecialCharTok{\textgreater{}} \DecValTok{1} \SpecialCharTok{\&}\NormalTok{ freq }\SpecialCharTok{\textgreater{}=} \DecValTok{2}\NormalTok{)}
\end{Highlighting}
\end{Shaded}

\begin{verbatim}
              keyword ngram freq
33           e---mail     3    9
34             --mail     2    9
49   customer-service     2    5
57            new-one     2    4
72   first-generation     2    3
74           usb-port     2    3
77         list-price     2    3
80        apple-store     2    3
83      apple-product     2    3
93      serial-number     2    3
95       tech-support     2    3
97           web-site     2    3
104          year-old     2    2
105         many-site     2    2
107        long-story     2    2
109          next-day     2    2
110   wifi-connection     2    2
111         long-time     2    2
112         first-one     2    2
114        wifi-issue     2    2
116        same-thing     2    2
118     free-shipping     2    2
123  software-problem     2    2
124       credit-card     2    2
125 electronic-device     2    2
128      day-warranty     2    2
\end{verbatim}

\begin{Shaded}
\begin{Highlighting}[]
\CommentTok{\# Extract keywords for Samsung (nouns/adjectives)}
\NormalTok{tr\_kw\_s }\OtherTok{\textless{}{-}} \FunctionTok{textrank\_keywords}\NormalTok{(sent\_df\_smsng}\SpecialCharTok{$}\NormalTok{lemma, }
                             \AttributeTok{relevant =}\NormalTok{ sent\_df\_smsng}\SpecialCharTok{$}\NormalTok{pos }\SpecialCharTok{\%in\%} \FunctionTok{c}\NormalTok{(}\StringTok{"NOUN"}\NormalTok{, }\StringTok{"ADJ"}\NormalTok{))}
\CommentTok{\# Filter: multi{-}word, freq \textgreater{}= 4}
\FunctionTok{subset}\NormalTok{(tr\_kw\_s}\SpecialCharTok{$}\NormalTok{keywords, ngram }\SpecialCharTok{\textgreater{}} \DecValTok{1} \SpecialCharTok{\&}\NormalTok{ freq }\SpecialCharTok{\textgreater{}=} \DecValTok{4}\NormalTok{)}
\end{Highlighting}
\end{Shaded}

\begin{verbatim}
        keyword ngram freq
62 same-problem     2    4
65 model-number     2    4
66 touch-screen     2    4
67     web-page     2    4
80    same-loop     2    4
\end{verbatim}

We can do the same procedure for sentences, but this time for
both brands, applying the textrank function to them. We can also
plot the ranking of one of the brands to see this visually. The
ranking procedure might take a long time for larger data sets, but there
are dimensionality reduction algorithms which are useful in such cases,
for example, Minhash (Broder et al. 1998).

\begin{Shaded}
\begin{Highlighting}[]
\CommentTok{\# Extract unique sentences for Apple and Samsung}
\NormalTok{sentences\_apple }\OtherTok{\textless{}{-}} \FunctionTok{unique}\NormalTok{(sent\_df\_apple[, }\FunctionTok{c}\NormalTok{(}\StringTok{"textrank\_id"}\NormalTok{, }\StringTok{"sentence"}\NormalTok{)])}
\NormalTok{sentences\_smsng }\OtherTok{\textless{}{-}} \FunctionTok{unique}\NormalTok{(sent\_df\_smsng[, }\FunctionTok{c}\NormalTok{(}\StringTok{"textrank\_id"}\NormalTok{, }\StringTok{"sentence"}\NormalTok{)])}

\CommentTok{\# Filter terms by POS tags for noun and adjective, select relevant columns}
\NormalTok{term\_apple }\OtherTok{\textless{}{-}}\NormalTok{ sent\_df\_apple[sent\_df\_apple}\SpecialCharTok{$}\NormalTok{pos }\SpecialCharTok{\%in\%} \FunctionTok{c}\NormalTok{(}\StringTok{"NOUN"}\NormalTok{, }\StringTok{"ADJ"}\NormalTok{), }
                            \FunctionTok{c}\NormalTok{(}\StringTok{"textrank\_id"}\NormalTok{, }\StringTok{"lemma"}\NormalTok{)]}
\NormalTok{term\_smsng }\OtherTok{\textless{}{-}}\NormalTok{ sent\_df\_smsng[sent\_df\_smsng}\SpecialCharTok{$}\NormalTok{pos }\SpecialCharTok{\%in\%} \FunctionTok{c}\NormalTok{(}\StringTok{"NOUN"}\NormalTok{, }\StringTok{"ADJ"}\NormalTok{), }
                            \FunctionTok{c}\NormalTok{(}\StringTok{"textrank\_id"}\NormalTok{, }\StringTok{"lemma"}\NormalTok{)]}

\CommentTok{\# Apply textrank to find relevant sentences}
\NormalTok{tr\_apple }\OtherTok{\textless{}{-}} \FunctionTok{textrank\_sentences}\NormalTok{(}\AttributeTok{data =}\NormalTok{ sentences\_apple, }\AttributeTok{terminology =}\NormalTok{ term\_apple)}
\NormalTok{tr\_smsng }\OtherTok{\textless{}{-}} \FunctionTok{textrank\_sentences}\NormalTok{(}\AttributeTok{data =}\NormalTok{ sentences\_smsng, }\AttributeTok{terminology =}\NormalTok{ term\_smsng)}

\CommentTok{\# Plot pagerank for Apple\textquotesingle{}s textrank results}
\FunctionTok{plot}\NormalTok{(}\FunctionTok{sort}\NormalTok{(tr\_apple}\SpecialCharTok{$}\NormalTok{pagerank}\SpecialCharTok{$}\NormalTok{vector, }\AttributeTok{decreasing =} \ConstantTok{TRUE}\NormalTok{), }\AttributeTok{type =} \StringTok{"b"}\NormalTok{, }
     \AttributeTok{ylab =} \StringTok{"Pagerank"}\NormalTok{, }\AttributeTok{main =} \StringTok{"Textrank"}\NormalTok{)}
\end{Highlighting}
\end{Shaded}

\includegraphics{9.Chapter_new_files/figure-pdf/sum5a-1.pdf}

We have extracted the top three sentences for Apple, although we could have
extracted as many as we wanted. They appear in ranking order.

\begin{Shaded}
\begin{Highlighting}[]
\CommentTok{\# Top 3 sentences for Apple}
\NormalTok{sum\_apple }\OtherTok{\textless{}{-}} \FunctionTok{summary}\NormalTok{(tr\_apple, }\AttributeTok{n =} \DecValTok{3}\NormalTok{, }\AttributeTok{keep.sentence.order =} \ConstantTok{TRUE}\NormalTok{)}
\FunctionTok{print}\NormalTok{(sum\_apple)}
\end{Highlighting}
\end{Shaded}

\begin{verbatim}
[1] "Got the product and had problems with it right away."                                
[2] "Just after I ordered and received this product, Apple released their newest IPad II."
[3] "Nor can I review all the items I purchased to use with the product."                 
\end{verbatim}

Here we have extracted the top three sentences for Samsung.

\begin{Shaded}
\begin{Highlighting}[]
\CommentTok{\# Top 3 sentences for Samsung}
\NormalTok{sum\_smsng }\OtherTok{\textless{}{-}} \FunctionTok{summary}\NormalTok{(tr\_smsng, }\AttributeTok{n =} \DecValTok{3}\NormalTok{, }\AttributeTok{keep.sentence.order =} \ConstantTok{TRUE}\NormalTok{)}
\FunctionTok{print}\NormalTok{(sum\_smsng)}
\end{Highlighting}
\end{Shaded}

\begin{verbatim}
[1] "Problem solved with the tablet."     
[2] "At first the tablet works just fine."
[3] "Love this tablet."                   
\end{verbatim}

Interestingly, sentences with \texttt{textrank} are shorter than in the
results obtained using the previous package. We can see that a comparison with other brands is made very often, especially in negative
reviews. As an exercise, plot pagerank for Samsung's textrank results

\subsection{A note}\label{a-summarization-note}

Abstractive summarizing is a more complex form
compared to extractive summarizing. It involves generating new
sentences that convey the main ideas of the original text. In R,
abstractive summarizing is not straightforward due to the lack of direct support for complex neural-network models within R itself. However, you can use R to interface with
other tools or libraries which support abstractive summarizing, such as
Python libraries via reticulate or external APIs. But in the next
chapter, we will discuss more advanced tools for carrying out text
analysis, and these will allow us to generate text, too.

\subsection{Conclusion and exercises}\label{conclusions-and-exercises}

Exercise 1: Select all the entities which refer to a person's name. Are
they related to the data set?\\
Exercise 2: Visualize a series of nouns for two brands over a
timeline.\\
Exercise 3: Create a model of linguistic complexity over a timeline.\\
Exercise 4: Read \texttt{spacyr}'s documentation and use a
consolidation function to merge multi-word entities.\\
Exercise 5: Visually plot the relative importance of sentences in the
\texttt{tr\_smsng} variable.

\emph{Conclusion starts}

The chapter began with an introduction to named entity recognition, its
definition, and why this method is useful in marketing. We discussed several packages used for
NER, with \texttt{spacyr} being the package of
choice for performing analysis. The package is discussed in detail,
and its types of annotation are also presented. 
We have focused on the extraction of entities, along with other parts of speech. We discussed
the information extraction based on specific POS, such as
adjectives whose frequency is also presented as a timeline on a brand-comparison plot. Continuing with visual representation, we used this to examine parts of speech dependencies, as well as entity,
nouns, and highlighting adjectives. A special mention was given to the
text complexity, which can offer insights into how customers express
their frustration in negative reviews. Summarizing is related to NER,
especially via the keywords which are parts-of-speech. We also introduced three R packages
which perform extractive summarizing. This chapter has
provided an overview of NER and summarizing, and complements the methods of
text analysis discussed in previous chapters.
